\section{Practical implications}
\label{sec:guide}

These suggestions come from four judges on one corpus. Each is cheap to check with one's own judge, since the harness
runs a judge through the six contexts in about 1,200 calls.

An evaluation in which the judge receives a correct record says little about the judge: with a correct record our
judges did what a lookup does. The more useful test removes links from the record and compares each judge with
itself on the affected items. If its catches there fall to the lookup's level, the judge contributes nothing beyond
the record on those items, and the completeness of the record becomes the real limit of the pipeline. That trade-off
can still be worth accepting, since the same record also spared the judges false alarms on decoys.

How the record is presented can matter, and not in the same way for every judge. For three of our judges a true
warning that the record was incomplete restored most catches on its gaps, but it also brought back the false alarms on
decoys that the record had removed, and the fourth judge recovered few. A builder considering such a warning should
test it on their own judge and items, count false alarms as well as catches, and make sure the warning is true.

Scoring should use reasons. Any-finding scoring inflates false alarms on correct answers and cannot distinguish a
judge that raises an irrelevant objection from one that says nothing. Where the record lists both replacements and
non-replacing updates, as the RFC index does, date controls and decoys come at no extra cost.

If the replacements in a target collection tend to look unlike the documents they replace, catches on such pairs
should be measured separately, since our judges were weakest there.
